\begin{pastexam}[label={px:ntu-111-dsaa-09}]{臺大資工所 111 學年度 第 9 題}
（單選）

(5 points) Consider an array of 8 elements being sorted using quicksort. It has just finished the
first pass of partitioning and pivot swapping, thus changing the original array into the
following array:
\[[7,11,16,10,17,1,18,30]\]
How many elements could have been the pivot?
\begin{enumerate}[label=(\Alph*)]
\item 1
\item 2
\item 3
\item 4
\item none of the other choices
\end{enumerate}
\end{pastexam}

\begin{pastsoln}
答案是 (B) 2。

第一次分割、把基準換到定位以後（資結書第~14.5~節），基準左邊的數都不比它大，右邊的數都不比它
小。 所以一個數「可能是基準」的條件是：左邊每個數都比它小，右邊每個數都比它大。 逐個檢查
（八個數都不同）：
\begin{center}
\begin{tabular}{cllc}
\toprule
數 & 左邊有比它大的？ & 右邊有比它小的？ & 可能是基準 \\
\midrule
7 & 無 & 有（1） & 否 \\
11 & 無 & 有（10、1） & 否 \\
16 & 無 & 有（10、1） & 否 \\
10 & 有（11、16） & 有（1） & 否 \\
17 & 無 & 有（1） & 否 \\
1 & 有（全部） & 無 & 否 \\
18 & 無 & 無 & 是 \\
30 & 無 & 無（右邊沒有數） & 是 \\
\bottomrule
\end{tabular}
\end{center}
快一點的做法：從左往右算「到目前為止的最大值」，從右往左算「到目前為止的最小值」，一個數同時
等於左邊的前綴最大值、右邊的後綴最小值，就可能是基準。 前綴最大值是 $7, 11, 16, 16, 17, 17,
18, 30$，後綴最小值是 $1, 1, 1, 1, 1, 1, 18, 30$，只有 $18$ 和 $30$ 兩個位置都相等。

只有 $18$、$30$ 兩個數可能是基準，選 (B)。 (A) 漏了一個；(C)、(D) 把不符合條件的數也算進去。
\end{pastsoln}
